% ============================================================
% Digitalización de señales y aliasing (LAB1)
%
% Construido a partir de la implementación real del proyecto:
%   - firmware/part1_digitization/         -> muestreo periódico, 500 Hz, ADC 12 bits
%   - firmware/part2_dac_reconstruction/    -> reconstrucción vía DAC
%   - firmware/part2_pwm_reconstruction/    -> reconstrucción vía PWM + filtro
%   - firmware/part3_bme280_sensor/         -> digitalización BME280 (I2C), 100 Hz
%   - firmware/part3_encoder_motor/         -> digitalización encoder/motor, 200 Hz
%   - pc_logger/serial_logger.py            -> captura eficiente a CSV, sin graficar
%   - matlab/*.m                            -> FFT, estadística, tablas de resultados
%
% Ningún hardware físico (tarjeta ESP32-WROOM, osciloscopio, generador de
% señales, BME280, encoder/motoreductor) estuvo disponible en este entorno;
% las tablas de resultados se marcan explícitamente como pendientes de
% captura real en lugar de contener datos inventados.
%
% Para compilar: pdflatex main.tex (dos pasadas si se usa \ref/\eqref)
% ============================================================

\documentclass[conference]{IEEEtran}

% ============================================================
% PAQUETES
% ============================================================

\usepackage[utf8]{inputenc}
\usepackage[T1]{fontenc}
\usepackage[spanish]{babel}

\usepackage{amsmath}
\usepackage{amssymb}

\usepackage{graphicx}
\usepackage{float}
\usepackage{booktabs}
\usepackage{array}
\usepackage{listings}

\usepackage{cite}
\usepackage{url}

\usepackage{xcolor}

\graphicspath{{figuras/}}

\lstset{
  basicstyle=\ttfamily\scriptsize,
  breaklines=true,
  frame=single,
  columns=fullflexible
}

% ============================================================
% INFORMACIÓN DEL DOCUMENTO
% ============================================================

\title{
Digitalización de señales analógicas periódicas, señales de sensores, reconstrucción y aliasing sobre un microcontrolador de 32 bits
}

\author{
\IEEEauthorblockN{Daniel García Araque}
\IEEEauthorblockA{
Programa de Ingeniería Mecatrónica\\
Universidad Militar Nueva Granada\\
Bogotá, Colombia\\
Código 7004185\\
est.daniel.garciaa@unimilitar.edu.co
}
\and
\IEEEauthorblockN{Miguel Ángel Gordillo Mejía}
\IEEEauthorblockA{
Programa de Ingeniería Mecatrónica\\
Universidad Militar Nueva Granada\\
Bogotá, Colombia\\
Código 7004909\\
est.miguel.gordillo@unimilitar.edu.co
}
}

% ============================================================
% DOCUMENTO
% ============================================================

\begin{document}

\maketitle

% ============================================================
% RESUMEN
% ============================================================

\begin{abstract}

Este informe documenta el diseño e implementación de un sistema de digitalización de señales basado en un microcontrolador de 32 bits (ESP32-WROOM, vía el núcleo Arduino para ESP32; la guía de práctica nombra una tarjeta STM32 específicamente para la etapa de reconstrucción PWM, pero este proyecto se desarrolla sobre ESP32 al ser la tarjeta realmente disponible — ver Sección de Metodología para las diferencias de hardware relevantes, en particular la resolución del DAC integrado), desarrollado para la práctica \texttt{LAB1} del curso. El sistema abarca cuatro etapas: (1) digitalización de señales analógicas periódicas mediante interrupción por temporizador a $500$ muestras/segundo con conversión análogo-digital de $12$ bits; (2) reconstrucción de la señal digitalizada mediante el módulo DAC de la tarjeta y, alternativamente, mediante PWM a $5$ kHz seguido de un filtro pasa-bajos analógico externo; (3) digitalización de señales de sensores —un sensor de presión, temperatura y humedad BME280 por I2C, a $100$ muestras/segundo (límite impuesto por el propio sensor, no $200$ como en la sección de encoder — ver Sección de Metodología), y un encoder acoplado a un motorreductor DC a $200$ muestras/segundo—; y (4) identificación experimental del fenómeno de aliasing con frecuencias superiores a la de muestreo. La guía de práctica nombra originalmente una unidad de medición inercial (IMU) para esta tercera etapa; el sensor realmente disponible es un BME280 en su lugar, por lo que esta etapa mide presión/temperatura/humedad y no movimiento — ver Sección de Metodología para el detalle. Todo el firmware, el programa de captura eficiente hacia el computador (sin visualización en tiempo real, por requisito de la práctica) y los scripts de análisis en MATLAB (FFT, estadística descriptiva, generación de las tablas solicitadas) están implementados y documentados en este repositorio. \textbf{No se dispuso de hardware físico} (tarjeta ESP32-WROOM, osciloscopio, generador de señales, BME280 ni motorreductor) en el entorno de desarrollo de este informe; en consecuencia, las tablas de resultados se presentan con sus valores marcados explícitamente como pendientes de captura real, y las preguntas de reflexión de la guía se responden desde el fundamento teórico del muestreo, la cuantización y la reconstrucción de señales, en lugar de a partir de mediciones inventadas. Todo el código de este proyecto está disponible en el repositorio público en \url{https://github.com/DanielAraqueStudios/Signals-frecuency/tree/main/LAB_TESTING/LAB1}.

\end{abstract}

% ============================================================
% PALABRAS CLAVE
% ============================================================

\begin{IEEEkeywords}

Muestreo, cuantización, aliasing, conversión análogo-digital, conversión digital-análoga, PWM, filtro pasa-bajos, BME280, encoder, ESP32, Transformada Rápida de Fourier.

\end{IEEEkeywords}

% ============================================================
% I. INTRODUCCIÓN
% ============================================================

\section{Introducción}

La conversión de una señal física continua en una representación digital procesable es la base de todo sistema de adquisición de datos moderno. Este proceso involucra tres decisiones de diseño fundamentales: la frecuencia a la que se toman las muestras (muestreo), el número de niveles discretos con los que se representa cada muestra (cuantización), y el mecanismo mediante el cual esa representación digital puede recuperarse de vuelta a una señal analógica (reconstrucción). Cada una de estas decisiones introduce una fuente distinta de pérdida de información frente a la señal física original, y cada una es controlable una vez se comprende su mecanismo.

Esta práctica desarrolla un sistema de digitalización completo sobre un microcontrolador de 32 bits, cubriendo las tres etapas mencionadas más un cuarto tema —el fenómeno de aliasing— que surge precisamente cuando la primera de esas decisiones, la frecuencia de muestreo, no se elige adecuadamente respecto al contenido frecuencial de la señal de entrada. El sistema implementado (\texttt{firmware/}) captura señales periódicas y de sensores mediante interrupciones de temporizador de alta precisión, las transmite eficientemente al computador mediante un protocolo binario compacto (\texttt{pc\_logger/}), y sus resultados se analizan estadística y frecuencialmente en MATLAB (\texttt{matlab/}).

El objetivo es doble: por un lado, implementar y verificar en la práctica los mecanismos de muestreo, cuantización, transmisión eficiente de datos y reconstrucción de señales; por otro, caracterizar experimentalmente las fuentes reales de error de un sistema de digitalización —diferencia entre frecuencia normalizada calculada y observada, desfase entre una señal reconstruida y su original, ruido en sensores reales— en contraste con sus predicciones puramente teóricas.

% ============================================================
% II. FUNDAMENTO TEÓRICO
% ============================================================

\section{Fundamento teórico}

\subsection{Muestreo y frecuencia normalizada}
\label{sec:muestreo}

Una señal analógica $x(t)$ se convierte en una secuencia discreta mediante el muestreo periódico \cite{oppenheim}:

\begin{equation}
x[n] = x(nT_s), \qquad T_s = \frac{1}{f_s}
\end{equation}

La \textbf{frecuencia normalizada} de una componente de frecuencia $f$ dentro de una señal muestreada a $f_s$ se define como:

\begin{equation}
f_{\text{norm}} = \frac{f}{f_s} \quad \text{(ciclos/muestra)}
\label{eq:norm}
\end{equation}

y equivale al número de ciclos que dicha componente completa entre una muestra y la siguiente. De forma equivalente, el número de \textbf{muestras por ciclo} es el recíproco de la Ec.~\ref{eq:norm}: $N_{\text{muestras/ciclo}} = f_s/f$.

\subsection{Teorema de muestreo y aliasing}

El teorema de Nyquist establece que una señal se puede reconstruir sin ambigüedad únicamente si:

\begin{equation}
f_s \geq 2 f_{\max}
\end{equation}

Cuando esta condición no se cumple, ocurre \textbf{aliasing}: una componente de frecuencia $f > f_s/2$ se manifiesta en la señal muestreada como una componente de frecuencia \emph{aparente} distinta, obtenida al ``plegar'' $f$ hacia el intervalo $[0, f_s/2]$ mediante reflexiones sucesivas alrededor de los múltiplos de $f_s/2$:

\begin{equation}
f_{\text{aliased}} = \left| f - k f_s \right|, \quad k = \text{round}(f/f_s)
\label{eq:aliasing}
\end{equation}

para el $k$ entero que deja $f_{\text{aliased}} \in [0, f_s/2]$. Esta es la relación que implementa \texttt{matlab/analyze\_aliasing.m} para calcular la frecuencia normalizada teórica de la sección opcional de esta práctica.

\subsection{Cuantización}

La cuantización aproxima cada muestra continua en amplitud a uno de $2^b$ niveles discretos, donde $b$ es la resolución del convertidor —$12$ bits en este sistema, es decir $4096$ niveles—. El error de cuantización máximo es de medio nivel de cuantización (LSB/2), y su magnitud absoluta depende del rango de referencia del convertidor: para un ADC de $12$ bits con referencia de $3.3$ V, un LSB equivale a $3.3/4095 \approx 0.8$ mV.

\subsection{Transformada Rápida de Fourier}

La FFT calcula eficientemente la Transformada Discreta de Fourier \cite{lyons}, permitiendo identificar el contenido frecuencial de una señal ya muestreada. La componente en $0$ Hz del espectro corresponde al nivel de continua (DC) de la señal —su valor medio—, y su magnitud es directamente proporcional a la amplitud de dicho nivel DC. El pico correspondiente a una componente sinusoidal de frecuencia $f$ y amplitud $A$ aparece en la posición $f$ del espectro con una magnitud proporcional a $A$; con la normalización $2/N$ empleada en \texttt{matlab/analyze\_periodic\_signal.m} (multiplicar el espectro de un solo lado por $2$ y dividir por el número de muestras $N$), la magnitud del pico converge directamente al valor de $A$ para una señal sinusoidal pura sin ventaneo.

\subsection{Reconstrucción: DAC vs. PWM + filtro pasa-bajos}
\label{sec:reconstruccion}

La reconstrucción de una señal analógica a partir de sus muestras digitales puede realizarse mediante:

\begin{itemize}
\item \textbf{Conversión Digital-Análoga (DAC) directa}: cada muestra se convierte inmediatamente a un nivel de voltaje sostenido hasta la siguiente muestra (retenedor de orden cero, \emph{zero-order hold}). La salida resultante es una función escalonada cuyo número de escalones por ciclo es idéntico al número de muestras por ciclo capturadas —la misma cantidad calculada en la Sección~\ref{sec:muestreo}—, ya que cada muestra produce exactamente un escalón.
\item \textbf{PWM seguido de un filtro pasa-bajos analógico}: el valor de cada muestra se codifica como el ciclo útil (\emph{duty cycle}) de una señal PWM de alta frecuencia (aquí, $5$ kHz, diez veces la frecuencia de muestreo). El filtro pasa-bajos promedia esa señal de alta frecuencia, recuperando el valor medio de cada ciclo útil como un nivel de voltaje continuo por muestra —efectivamente un DAC construido con un promediador RC/activo en lugar de un convertidor dedicado.
\end{itemize}

Ambos métodos introducen un \textbf{desfase} (retardo) entre la señal original y la reconstruida, pero de naturaleza distinta: el DAC directo introduce principalmente el retardo de procesamiento discreto (el tiempo entre la captura de una muestra y su escritura en el DAC, del orden de microsegundos, y el propio retenedor de orden cero, que introduce medio periodo de muestreo de retardo de grupo en promedio); el filtro pasa-bajos del método PWM añade, además, su propio retardo de grupo, dependiente de su orden y frecuencia de corte, que típicamente domina sobre el retardo del retenedor de orden cero y por tanto no es igual al desfase del método DAC. Ninguno de los dos desfases es necesariamente constante con la frecuencia de la señal de entrada si el filtro no tiene fase lineal en la banda de interés —una característica de diseño de filtro, no del método de reconstrucción en sí.

% ============================================================
% III. METODOLOGÍA
% ============================================================

\section{Metodología}

El sistema se implementó sobre una tarjeta ESP32-WROOM mediante el núcleo Arduino para ESP32, en cinco sketches independientes (\texttt{firmware/}), un programa de captura en Python (\texttt{pc\_logger/serial\_logger.py}) y cuatro scripts de análisis en MATLAB (\texttt{matlab/}). La guía de práctica nombra una tarjeta STM32 específicamente para la etapa de reconstrucción PWM; este proyecto se desarrolla sobre ESP32 al ser la tarjeta realmente disponible, con una diferencia de hardware relevante frente a un STM32: el DAC integrado del ESP32 (\texttt{part2\_dac\_reconstruction.ino}) es de solo $8$ bits (pines \texttt{GPIO25}/\texttt{GPIO26}), no de $12$ bits; la captura ADC sigue siendo de $12$ bits según lo exigido, pero el código se reduce a $8$ bits (\texttt{code >> 4}) únicamente para la escritura al DAC. \textbf{Ningún componente de hardware} —tarjeta ESP32-WROOM, osciloscopio, generador de señales, IMU o encoder/motorreductor— estuvo disponible en el entorno de desarrollo de este informe; todo el código se escribió para compilar y ejecutarse tal como está, pero no ha sido flasheado ni probado en hardware real.

\subsection{Digitalización de señales periódicas}

\texttt{firmware/part1\_digitization/part1\_digitization.ino} configura una interrupción de temporizador de hardware (\texttt{hw\_timer\_t}, API \texttt{timerBegin}/\texttt{timerAlarmWrite}) cada $2$ ms ($500$ muestras/segundo). La rutina de interrupción alterna un pin de depuración (verificable con osciloscopio, para confirmar el periodo de $2$ ms con el error máximo del $5\,\%$ exigido) y captura una muestra del ADC configurado a $12$ bits; el envío por Serial se realiza en \texttt{loop()}, no dentro de la interrupción, dado que \texttt{Serial.write()} no está garantizado como seguro en contexto de interrupción sobre el núcleo Arduino de ESP32 (ver \texttt{firmware/README.md} para el patrón completo). El código crudo se transmite al PC en una trama binaria de $3$ bytes por muestra; la conversión a voltios se realiza en el lado del PC (\texttt{pc\_logger}), manteniendo la rutina de interrupción lo más corta posible.

\subsection{Transmisión eficiente de datos}

Se descartó $9600$ bit/s por insuficiente: a $500$ muestras/s con tramas de $3$ bytes se requieren $1500$ B/s, cifra que ya excede la capacidad útil de esa velocidad ($\approx 960$ B/s tras bits de inicio/parada). Se emplearon $115200$ baudios para las partes 1 y 2 ($\approx 11520$ B/s, margen de $7.7\times$) y $230400$ baudios para el sensor IMU ($\approx 23040$ B/s, margen de $8.2\times$ sobre sus $2800$ B/s requeridos), justificación completa en \texttt{firmware/README.md}. \texttt{pc\_logger/serial\_logger.py} implementa la recepción: sincroniza mediante un byte distintivo por tipo de trama, decodifica los canales a unidades físicas, y escribe directamente a CSV en disco —sin visualización ni graficación en tiempo real, por requisito explícito de la guía—. La lógica de sincronización y decodificación se valida con $14$ pruebas unitarias (\texttt{pc\_logger/tests/test\_serial\_logger.py}) contra flujos de bytes sintéticos, dado que no se dispuso de un puerto serie real para probarla de extremo a extremo.

\subsection{Reconstrucción de la señal digitalizada}

Dos sketches implementan los dos métodos de reconstrucción exigidos, ambos concurrentes con la misma interrupción de $2$ ms de captura:

\begin{itemize}
\item \texttt{part2\_dac\_reconstruction.ino}: el código ADC de $12$ bits capturado se reduce a los $8$ bits nativos del DAC del ESP32 (\texttt{code >> 4}) y se escribe inmediatamente, dentro de la misma rutina de interrupción.
\item \texttt{part2\_pwm\_reconstruction.ino}: mediante el periférico LEDC del ESP32, el código ADC se traduce a un ciclo útil PWM (resolución de $10$ bits) entre $1\,\%$ y $99\,\%$ (según lo recomendado en la guía), a una frecuencia PWM de $5$ kHz —simultáneamente diez veces la frecuencia de muestreo y superior al mínimo de $5$ kHz exigido—. La salida PWM requiere un filtro pasa-bajos analógico externo (no implementado en firmware) con ganancia $0$ dB en $500$ Hz y $\leq\!-20$ dB en $600$ Hz; dado que un filtro RC de un solo polo solo atenúa $-20$ dB/década, esta especificación exige un filtro activo de al menos segundo orden (por ejemplo, Sallen-Key), documentado como pendiente de diseño/construcción física.
\end{itemize}

\subsection{Digitalización de sensores}

\texttt{part3\_imu\_sensor.ino} muestrea una IMU de seis canales (acelerómetro + giroscopio) por I2C a $200$ muestras/segundo, asumiendo el mapa de registros de una MPU-6050 (ajustable si se usa otro modelo). La conversión a unidades físicas (m/s$^2$, rad/s) usa las sensibilidades del fabricante para el rango de escala completa por defecto ($\pm 2$g, $\pm 250\,^{\circ}$/s), implementada en \texttt{pc\_logger/serial\_logger.py::imu\_decode}.

\texttt{part3\_encoder\_motor.ino} implementa la estrategia de \emph{conteo de pulsos por intervalo} (opción 1 de la guía) sobre una ventana de $5$ ms, mediante una interrupción por flanco en el canal A del encoder y el signo determinado por el canal B. Esta estrategia se prefirió sobre la medición de tiempo entre pulsos porque es más robusta al rango de velocidades típico de un motorreductor con carga (velocidades moderadas y relativamente constantes), donde medir el tiempo entre pulsos individuales es más sensible al ruido de conmutación del encoder. El valor \texttt{PULSES\_PER\_REV} debe determinarse experimentalmente con un tacómetro antes de interpretar los resultados en unidades angulares, según lo exige la guía.

\subsection{Análisis en MATLAB}

Los cuatro scripts en \texttt{matlab/} leen los CSV producidos por \texttt{pc\_logger/serial\_logger.py} y calculan exactamente los valores solicitados en cada tabla de resultados de la Sección~\ref{sec:resultados}: muestras por ciclo y frecuencia normalizada calculada/observada (\texttt{analyze\_periodic\_signal.m}), estadística descriptiva y ancho de banda por canal IMU (\texttt{analyze\_imu.m}), estadística y ancho de banda de velocidad angular (\texttt{analyze\_encoder.m}), y frecuencia normalizada calculada (plegada, Ec.~\ref{eq:aliasing}) vs. observada para el análisis opcional de aliasing (\texttt{analyze\_aliasing.m}). Ninguno de los cuatro se ha ejecutado contra datos reales en este informe.

% ============================================================
% IV. RESULTADOS
% ============================================================

\section{Resultados}
\label{sec:resultados}

\textbf{Ningún resultado de esta sección proviene de hardware real}: no se dispuso de tarjeta STM32, osciloscopio, generador de señales, IMU ni encoder/motorreductor en el entorno de desarrollo de este informe. Las columnas calculables desde la teoría del muestreo (frecuencia normalizada calculada, muestras por ciclo teóricas) se completan a continuación; las columnas que requieren una captura real (frecuencia normalizada \emph{observada}, amplitudes/frecuencias reconstruidas, estadística de sensores reales) se marcan como \emph{pendiente} y deben completarse ejecutando \texttt{matlab/*.m} sobre los CSV generados por \texttt{pc\_logger/serial\_logger.py} una vez el hardware esté disponible.

\subsection{Digitalización de señales periódicas}

\begin{table}[!t]
\centering
\caption{Frecuencia normalizada — señal periódica ($f_s = 500$ Hz)}
\label{tab:periodica}
\small
\begin{tabular}{@{}r r r r r@{}}
\toprule
$f$ (Hz) & Muestras/ciclo & $f_{\text{norm}}$ calc. & $f_{\text{norm}}$ obs. & Diferencia \\
\midrule
50 & 10 & $1/10 = 0.100$ & pendiente & pendiente \\
75 & 6.667 & $3/20 = 0.150$ & pendiente & pendiente \\
100 & 5 & $1/5 = 0.200$ & pendiente & pendiente \\
125 & 4 & $1/4 = 0.250$ & pendiente & pendiente \\
\bottomrule
\end{tabular}
\end{table}

Las columnas de muestras/ciclo y frecuencia normalizada calculada de la Tabla~\ref{tab:periodica} se obtienen directamente de $f_s/f$ y $f/f_s$ (Ec.~\ref{eq:norm}), sin requerir captura alguna; nótese que $75$ Hz no produce un número entero de muestras por ciclo ($6.667$), a diferencia de $50$, $100$ y $125$ Hz, un punto retomado en la Discusión.

\subsection{Reconstrucción de la señal digitalizada}

\begin{table}[!t]
\centering
\caption{Comparación entrada vs. señales reconstruidas}
\label{tab:reconstruccion}
\small
\begin{tabular}{@{}r r r r r r@{}}
\toprule
$f_{\text{in}}$ (Hz) & $A_{\text{in}}$ (V) & $f_{\text{DAC}}$ (Hz) & $A_{\text{DAC}}$ (V) & $f_{\text{PWM}}$ (Hz) & $A_{\text{PWM}}$ (V) \\
\midrule
50 & pendiente & pendiente & pendiente & pendiente & pendiente \\
75 & pendiente & pendiente & pendiente & pendiente & pendiente \\
100 & pendiente & pendiente & pendiente & pendiente & pendiente \\
125 & pendiente & pendiente & pendiente & pendiente & pendiente \\
\bottomrule
\end{tabular}
\end{table}

Todas las columnas de la Tabla~\ref{tab:reconstruccion} requieren observación directa en el osciloscopio de las señales reconstruidas por \texttt{part2\_dac\_reconstruction.ino} y \texttt{part2\_pwm\_reconstruction.ino} (con su filtro externo), y quedan pendientes.

\subsection{Digitalización de sensores — IMU}

\begin{table}[!t]
\centering
\caption{Estadística IMU — sensor quieto}
\label{tab:imu-quieto}
\small
\begin{tabular}{@{}l r r r r@{}}
\toprule
Canal & Media & Desv. est. & Nivel DC & Ancho de banda \\
\midrule
$a_x$ & pendiente & pendiente & pendiente & pendiente \\
$a_y$ & pendiente & pendiente & pendiente & pendiente \\
$a_z$ & pendiente & pendiente & pendiente & pendiente \\
$g_x$ & pendiente & pendiente & pendiente & pendiente \\
$g_y$ & pendiente & pendiente & pendiente & pendiente \\
$g_z$ & pendiente & pendiente & pendiente & pendiente \\
\bottomrule
\end{tabular}
\end{table}

\begin{table}[!t]
\centering
\caption{Estadística IMU — sensor en movimiento}
\label{tab:imu-movimiento}
\small
\begin{tabular}{@{}l r r r r@{}}
\toprule
Canal & Media & Desv. est. & Nivel DC & Ancho de banda \\
\midrule
$a_x$ & pendiente & pendiente & pendiente & pendiente \\
$a_y$ & pendiente & pendiente & pendiente & pendiente \\
$a_z$ & pendiente & pendiente & pendiente & pendiente \\
$g_x$ & pendiente & pendiente & pendiente & pendiente \\
$g_y$ & pendiente & pendiente & pendiente & pendiente \\
$g_z$ & pendiente & pendiente & pendiente & pendiente \\
\bottomrule
\end{tabular}
\end{table}

Las Tablas~\ref{tab:imu-quieto} y \ref{tab:imu-movimiento} se completan ejecutando \texttt{matlab/analyze\_imu.m} sobre \texttt{data/imu\_still/*.csv} y \texttt{data/imu\_moving/*.csv} respectivamente, generados con \texttt{pc\_logger/serial\_logger.py --mode imu}.

\subsection{Digitalización de sensores — encoder y motorreductor}

\begin{table}[!t]
\centering
\caption{Velocidad angular vs. voltaje aplicado}
\label{tab:encoder}
\small
\begin{tabular}{@{}r r r r@{}}
\toprule
$V$ (V) & $\bar{\omega}$ (rad/s) & Desv. est. & Ancho de banda \\
\midrule
3 & pendiente & pendiente & pendiente \\
6 & pendiente & pendiente & pendiente \\
9 & pendiente & pendiente & pendiente \\
12 & pendiente & pendiente & pendiente \\
\bottomrule
\end{tabular}
\end{table}

Los voltajes de la Tabla~\ref{tab:encoder} son valores de referencia sugeridos, sujetos a la tensión nominal real del motorreductor disponible; se completan ejecutando \texttt{matlab/analyze\_encoder.m} sobre \texttt{data/encoder/*.csv}, con \texttt{PULSES\_PER\_REV} verificado por tacómetro según la Sección de Metodología.

\subsection{Fenómeno de aliasing (opcional)}

\begin{table}[!t]
\centering
\caption{Frecuencia normalizada — régimen de aliasing ($f_s = 500$ Hz)}
\label{tab:aliasing}
\small
\begin{tabular}{@{}r r r r@{}}
\toprule
$f$ (Hz) & $f_{\text{norm}}$ calc. & $f_{\text{norm}}$ obs. & $f_{\text{reconstruida}}$ (Hz) \\
\midrule
250 & 0.500 & pendiente & pendiente \\
500 & 0.000 & pendiente & pendiente \\
550 & 0.100 & pendiente & pendiente \\
575 & 0.150 & pendiente & pendiente \\
950 & 0.100 & pendiente & pendiente \\
1000 & 0.000 & pendiente & pendiente \\
1050 & 0.100 & pendiente & pendiente \\
1075 & 0.150 & pendiente & pendiente \\
2000 & 0.000 & pendiente & pendiente \\
2050 & 0.100 & pendiente & pendiente \\
2075 & 0.150 & pendiente & pendiente \\
\bottomrule
\end{tabular}
\end{table}

La columna de frecuencia normalizada calculada de la Tabla~\ref{tab:aliasing} se obtiene de la Ec.~\ref{eq:aliasing} (plegado teórico alrededor de los múltiplos de $f_s/2 = 250$ Hz), sin requerir captura alguna. Nótese que $550$, $950$, $1050$, $2050$ Hz comparten la misma frecuencia normalizada calculada ($0.100$) que su contraparte de la Tabla~\ref{tab:periodica} no aliasada más cercana no aplica directamente, pero sí coinciden entre sí — el mismo patrón se repite para $575$/$1075$/$2075$ Hz (todas en $0.150$) y para $500$/$1000$/$2000$ Hz (todas en $0.000$, es decir, aparecen como nivel DC). Esta coincidencia es precisamente la firma del aliasing y se discute en la Sección~\ref{sec:discusion-aliasing}.

% ============================================================
% V. DISCUSIÓN
% ============================================================

\section{Discusión}

Esta sección responde las preguntas de reflexión planteadas por la guía de práctica. Al no disponerse de datos reales (Sección~\ref{sec:resultados}), las respuestas se fundamentan en el comportamiento teórico esperado del sistema implementado, señalando explícitamente qué debería verificarse una vez existan capturas reales.

\subsection{Digitalización de señales periódicas}

\textbf{Causas de las diferencias entre frecuencia normalizada calculada y observada.} Aun con una implementación correcta, se esperan diferencias no nulas por varias razones independientes: (1) la resolución en frecuencia de la FFT es $f_s/N$, por lo que el pico observado solo puede ubicarse en un múltiplo discreto de esa resolución, nunca exactamente en el valor teórico salvo coincidencia; (2) el error máximo permitido del $5\,\%$ en el periodo de la interrupción de temporizador (Sección de Metodología) se traduce directamente en un error proporcional de la frecuencia de muestreo efectiva $f_s$ y, por tanto, de toda frecuencia normalizada calculada con base en el valor nominal de $500$ Hz; (3) el generador de ondas de laboratorio tiene su propia tolerancia de frecuencia, no infinitamente precisa; y (4) cuando la ventana de captura no contiene un número entero de ciclos de la señal —como ocurre con $75$ Hz, que produce $6.667$ muestras/ciclo, no entero, según la Tabla~\ref{tab:periodica}—, se produce \emph{fuga espectral} (\emph{spectral leakage}): la energía de la componente sinusoidal se dispersa entre varios bins de la FFT en lugar de concentrarse en uno solo, ensanchando el pico observado y desplazando su máximo respecto a la frecuencia real.

\textbf{¿Están presentes en todas las filas?} Se espera que sí, en algún grado, dado que las causas (1)–(3) son inherentes al sistema y no dependen de la frecuencia específica. Sin embargo, la magnitud de la diferencia no debería ser uniforme: se espera que sea mayor para $75$ Hz que para $50$, $100$ y $125$ Hz, precisamente por la fuga espectral adicional de la causa (4), ya que estas tres últimas sí producen un número entero de muestras por ciclo dentro de una ventana de captura cuya duración sea un múltiplo entero del periodo de la señal.

\subsection{Análisis frecuencial}

El pico en $0$ Hz del espectro corresponde al nivel de continua (DC) de la señal, y su altura es directamente proporcional a la magnitud de dicho nivel DC —de hecho, tras la normalización $2/N$ empleada (Sección~\ref{sec:muestreo}), la altura del pico en $0$ Hz converge exactamente al valor del nivel DC configurado en el generador. La ubicación del pico correspondiente a la componente sinusoidal coincide con la frecuencia de dicha componente, y su amplitud converge al valor pico de la sinusoide (o a la mitad de su amplitud pico-a-pico) bajo la misma normalización. La formulación de ajuste sugerida es exactamente la implementada en \texttt{matlab/analyze\_periodic\_signal.m}: graficar $\text{mag} = \frac{2}{N}\left|X[k]\right|$ contra $\text{freq} = k \cdot f_s/N$ para $k = 0,\dots,N/2$, de modo que el eje horizontal quede directamente en Hz y el eje vertical directamente en las unidades físicas de la señal (voltios), sin factores de escala adicionales que el lector deba inferir.

\subsection{Reconstrucción de la señal digitalizada}

\textbf{Relación entre muestras por ciclo y escalones de la reconstrucción DAC.} Como se explica en la Sección~\ref{sec:reconstruccion}, cada muestra capturada produce exactamente un escalón en la salida del DAC (retenedor de orden cero), de modo que el número de escalones visibles por ciclo en la señal reconstruida es, por construcción, idéntico al número de muestras por ciclo calculado en la Tabla~\ref{tab:periodica} para esa misma frecuencia —no es una coincidencia experimental, sino una consecuencia directa de que el DAC solo puede cambiar de valor una vez por interrupción de muestreo.

\textbf{Causas del desfase de la reconstrucción PWM + filtro.} El desfase observado se debe principalmente al retardo de grupo introducido por el filtro pasa-bajos analógico externo, que depende de su orden y de su frecuencia de corte respecto a la componente de la señal; a este se suma un retardo menor y común a ambos métodos, el del propio retenedor de orden cero del muestreo (en promedio, medio periodo de muestreo). \textbf{¿Son iguales los desfases de PWM y DAC?} No necesariamente: el método DAC solo hereda el retardo del retenedor de orden cero, mientras que el método PWM añade, además, el retardo de grupo del filtro externo, que típicamente lo domina y por tanto produce un desfase total distinto (y usualmente mayor) que el del DAC directo. \textbf{¿Son los desfases constantes con la frecuencia de entrada?} El retardo del retenedor de orden cero es aproximadamente constante en muestras (medio periodo de muestreo), por lo que en tiempo absoluto no depende de la frecuencia de la señal de entrada; el retardo de grupo del filtro analógico, en cambio, generalmente sí varía con la frecuencia salvo que el filtro se diseñe específicamente con fase lineal en la banda de interés — por lo que el desfase del método PWM no debería considerarse constante en general.

\textbf{Ventajas, desventajas y aplicaciones.} El método DAC directo es más simple de implementar en firmware, no requiere circuitería analógica adicional, y su desfase es más predecible (depende solo del muestreo, no de un filtro externo); su desventaja es que su salida es intrínsecamente escalonada, con contenido espectral de alta frecuencia (armónicos del muestreo) que puede requerir su propio filtrado si se necesita una salida verdaderamente suave. El método PWM + filtro produce, tras un filtro bien diseñado, una salida más suave sin depender de la resolución del DAC de la tarjeta (útil si esta no tiene DAC dedicado, o si se necesita mayor potencia de salida mediante una etapa de conmutación); su desventaja es la complejidad y el costo adicional del filtro analógico externo, su retardo de grupo generalmente mayor y dependiente de la frecuencia, y la necesidad de un diseño de filtro cuidadoso para lograr suficiente atenuación de la portadora PWM sin distorsionar la señal de interés. En términos de aplicación, el DAC directo es preferible cuando se requiere baja latencia y simplicidad (por ejemplo, generación de referencias de control), mientras que PWM + filtro es preferible en escenarios de potencia (control de motores, actuadores) donde de todas formas se dispone de una etapa de conmutación y el filtrado forma parte natural del sistema.

\subsection{Digitalización de sensores — BME280}

La guía de práctica original formula esta subsección en torno a una IMU; al usarse un BME280 (presión, temperatura, humedad) en su lugar, no existen canales de aceleración ni de velocidad angular, y la comparación ``quieto'' vs. ``en movimiento'' no aplica directamente a un sensor sin partes ni ejes inerciales. Las preguntas se responden a continuación adaptadas a la comparación real implementada: condiciones \textbf{base} (ambiente estable) vs. \textbf{perturbadas} (Sección~\ref{sec:resultados}).

\textbf{Diferencias entre condiciones base y perturbadas.} En condiciones base, con el sensor en un ambiente térmicamente estable y sin fuentes cercanas de calor, humedad o presión diferencial, se espera que los tres canales muestren una media aproximadamente constante (la temperatura y humedad ambientales, y la presión atmosférica local) con una desviación estándar pequeña, atribuible principalmente al ruido electrónico interno del sensor y a la resolución efectiva de las fórmulas de compensación. Bajo perturbación deliberada (por ejemplo, una fuente de calor o el aliento acercados al sensor, o cubrirlo parcialmente), se espera un cambio significativo en la media de temperatura y humedad —ambas suben con el aliento o una fuente de calor cercana—, y un cambio más sutil, o incluso nulo, en la presión, ya que esta responde principalmente a la altitud/condiciones atmosféricas y no a perturbaciones locales pequeñas salvo que se module directamente el volumen de aire alrededor del sensor (por ejemplo, cubriéndolo).

\textbf{Identificación y características del ruido.} En condiciones base, el ruido de los tres canales es de banda ancha y de amplitud pequeña relativa al nivel DC dominante (temperatura/humedad/presión ambientales), consistente con ruido térmico/electrónico propio del sensor y con el redondeo introducido por la compensación en punto flotante —su ancho de banda calculado por \texttt{matlab/analyze\_bme280.m} debería ser relativamente amplio pero de magnitud pequeña. Bajo perturbación, el cambio dominante es un transitorio de baja frecuencia (el tiempo de establecimiento térmico/higrométrico del propio sensor, típicamente de segundos, mucho más lento que los $100$ Hz de muestreo) superpuesto al mismo ruido de fondo de banda ancha; a diferencia de un sensor inercial, aquí no se espera contenido de alta frecuencia genuino asociado a la perturbación, ya que las magnitudes físicas medidas (temperatura, presión, humedad) son intrínsecamente de dinámica lenta.

\textbf{Estrategias de mitigación de ruido.} Las estrategias aplicables incluyen filtrado digital pasa-bajos (un filtro de media móvil es especialmente adecuado, dada la dinámica lenta de las tres magnitudes), promediado de múltiples muestras —tolerable aquí sin penalizar la respuesta a eventos de interés, ya que estos ocurren en la escala de segundos, no de milisegundos—, aumento del sobremuestreo interno del propio BME280 (\texttt{ctrl\_meas}/\texttt{ctrl\_hum}, a costa de una tasa de muestreo aún menor que los $100$ Hz ya limitados por el sensor), y calibración de \emph{offset} en condiciones ambientales conocidas antes de cada sesión de medición si se requiere precisión absoluta, no solo relativa.

\subsection{Digitalización de sensores — encoder y motorreductor}

\textbf{Relación velocidad/voltaje aplicado.} Para un motor DC en estado estacionario y bajo carga aproximadamente constante, la relación entre la velocidad angular y el voltaje aplicado es aproximadamente lineal, consistente con el modelo de estado estable de un motor DC $\omega \approx (V - I R_a)/K_e$, donde para una carga fija la caída resistiva $I R_a$ varía poco frente a $V$, dejando una relación cercana a $\omega \propto V$ con un posible desplazamiento (voltaje mínimo de arranque, por fricción estática) antes de que el motor comience a girar.

\textbf{Características del ruido.} La señal de velocidad angular derivada del conteo de pulsos por intervalo tiene una fuente de ruido intrínseca de cuantización: el número de pulsos contados en cada ventana de $5$ ms es un entero, por lo que la velocidad estimada tiene una resolución mínima de $\frac{2\pi}{\texttt{PULSES\_PER\_REV} \times 0.005\,\text{s}}$ rad/s, que se manifiesta como ruido de cuantización de alta frecuencia superpuesto a la velocidad real, además de la variación genuina de velocidad instantánea del motor bajo carga.

\textbf{Estrategias de reducción de ruido.} Aumentar el tamaño de la ventana de conteo reduce el ruido de cuantización relativo (a costa de menor resolución temporal), promediar velocidades de varias ventanas consecutivas, o cambiar a la estrategia de medición de tiempo entre pulsos a velocidades bajas (donde produce menor ruido relativo que el conteo por intervalo) son las estrategias más directas; un filtro pasa-bajos digital adicional sobre la señal de velocidad estimada también es aplicable si la dinámica de interés es mucho más lenta que la resolución del encoder lo permite.

\subsection{Fenómeno de aliasing}
\label{sec:discusion-aliasing}

\textbf{¿Qué ocurre con los datos capturados cuando $f \geq 250$ Hz?} Al superar la frecuencia de Nyquist ($f_s/2 = 250$ Hz), la señal capturada deja de representar fielmente la frecuencia real del generador: en su lugar, el sistema captura una señal cuya frecuencia normalizada observada coincide con la frecuencia \emph{plegada} teórica de la Tabla~\ref{tab:aliasing} (Ec.~\ref{eq:aliasing}), indistinguible —a partir únicamente de los datos digitalizados— de una señal genuina de esa frecuencia más baja. Es exactamente lo que predice la columna de frecuencia normalizada calculada de la Tabla~\ref{tab:aliasing}: $550$, $950$, $1050$ y $2050$ Hz deberían observarse todos como $0.100$ ciclos/muestra ($50$ Hz aparentes); $575$, $1075$ y $2075$ Hz como $0.150$ ($75$ Hz aparentes); y $500$, $1000$ y $2000$ Hz como $0.000$, es decir, indistinguibles de una señal de continua (DC) pura.

\textbf{Conclusión sobre el proceso de digitalización.} La coincidencia teórica de frecuencias normalizadas entre múltiplos de $f_s$ separados exactamente $500$ Hz confirma que el proceso de muestreo es fundamentalmente periódico en el dominio de la frecuencia: cualquier frecuencia analógica y sus reflejos alrededor de los múltiplos de $f_s/2$ son indistinguibles una vez digitalizados, sin importar cuán distintas sean antes de muestrear.

\textbf{Explicación del fenómeno de aliasing a partir de los resultados.} El aliasing ocurre porque el proceso de muestreo, al tomar solo muestras discretas y periódicas de la señal, no conserva suficiente información para distinguir entre una componente de frecuencia $f$ y su reflejo $|f - kf_s|$ dentro de la banda $[0, f_s/2]$: ambas producen exactamente la misma secuencia de muestras. La reconstrucción posterior (sea DAC o PWM+filtro) no puede recuperar cuál de las dos frecuencias originó realmente los datos, y por tanto reproduce la frecuencia plegada, no la real —de ahí que la señal reconstruida en el osciloscopio para $f \geq 250$ Hz debería mostrar la frecuencia aparente de la Tabla~\ref{tab:aliasing}, no la frecuencia realmente configurada en el generador.

\textbf{Máxima frecuencia digitalizable sin afectar la información.} De acuerdo con el teorema de Nyquist, la máxima frecuencia que puede digitalizarse sin ambigüedad con $f_s = 500$ Hz es $f_s/2 = 250$ Hz —el propio límite en el que la Tabla~\ref{tab:aliasing} comienza su primera fila—, ya que cualquier frecuencia por encima de ese límite queda indistinguible de alguna frecuencia dentro de $[0, 250]$ Hz una vez muestreada.

% ============================================================
% VI. CONCLUSIONES
% ============================================================

\section{Conclusiones}

Se implementó un sistema completo de digitalización de señales sobre un microcontrolador de 32 bits, cubriendo muestreo periódico a $500$ Hz con $12$ bits de resolución, transmisión eficiente de datos ($115200$ baudios, justificados numéricamente frente al límite insuficiente de $9600$ bit/s), dos métodos de reconstrucción (DAC directo y PWM + filtro pasa-bajos), y digitalización de dos tipos de sensores por I2C: un BME280 (presión, temperatura, humedad) a $100$ Hz —límite propio del sensor, no $200$ Hz, ya que la guía original nombra una IMU en su lugar— y un encoder de motorreductor a $200$ Hz, además de análisis en MATLAB para las cuatro tablas de resultados solicitadas por la guía, incluyendo la sección opcional de aliasing.

El análisis teórico de las tablas de resultados —posible sin hardware físico para las columnas de frecuencia normalizada calculada— mostró que $75$ Hz es la única de las cuatro frecuencias de la sección de digitalización periódica que no produce un número entero de muestras por ciclo ($6.667$), anticipando una mayor fuga espectral y, por tanto, una mayor diferencia esperada entre frecuencia normalizada calculada y observada para esa fila específica frente a $50$, $100$ y $125$ Hz. De igual forma, el análisis teórico de la sección de aliasing predijo que $550$, $950$, $1050$ y $2050$ Hz deberían observarse todos con la misma frecuencia normalizada ($0.100$, equivalente a $50$ Hz aparentes), al igual que $575$/$1075$/$2075$ Hz ($0.150$) y $500$/$1000$/$2000$ Hz ($0.000$, indistinguibles de continua) — la firma característica del plegamiento espectral que define el aliasing, y una predicción verificable directamente una vez existan capturas reales.

Ningún resultado de esta práctica proviene de mediciones físicas: no se dispuso de tarjeta ESP32-WROOM, osciloscopio, generador de señales, BME280 ni encoder/motorreductor en el entorno de desarrollo de este informe. El firmware, el programa de captura eficiente, y los scripts de análisis en MATLAB están completos, documentados y listos para ejecutarse; el trabajo pendiente directo es flashear los cinco sketches en una tarjeta ESP32-WROOM real, capturar los datos con el hardware de laboratorio (osciloscopio, generador de señales, BME280 y motorreductor con encoder), y reemplazar cada valor marcado como \emph{pendiente} en las Tablas~\ref{tab:periodica}–\ref{tab:aliasing} por su medición real, verificando en el proceso las diferencias predichas teóricamente en esta sección de Discusión.

% ============================================================
% REFERENCIAS
% ============================================================

\begin{thebibliography}{00}

\bibitem{oppenheim}
A. V. Oppenheim, A. S. Willsky y S. H. Nawab,
\textit{Signals and Systems},
2nd ed., Prentice Hall, 1996.

\bibitem{lyons}
R. G. Lyons,
\textit{Understanding Digital Signal Processing},
3rd ed., Pearson, 2011.

\bibitem{esp32}
Espressif Systems,
\textit{ESP32 Technical Reference Manual},
2023. [En línea]. Disponible: \url{https://www.espressif.com/en/support/documents/technical-documents}

\bibitem{mpu6050}
InvenSense,
\textit{MPU-6000/MPU-6050 Register Map and Descriptions},
2013. [En línea]. Disponible: \url{https://invensense.tdi.com/wp-content/uploads/2015/02/MPU-6000-Register-Map1.pdf}

\bibitem{repo}
D. Araque,
\textit{Signals-frecuency} (repositorio de código),
2026. [En línea]. Disponible: \url{https://github.com/DanielAraqueStudios/Signals-frecuency/tree/main/LAB_TESTING/LAB1}

\end{thebibliography}


\end{document}
